\documentclass{article}

% Compiles with the ICML 2026 style package; remove or replace
% \usepackage{icml2026} if compiling without it.
\usepackage{microtype}
\usepackage{graphicx}
\usepackage{subcaption}
\usepackage{booktabs}
\usepackage{amsmath}
\usepackage{amssymb}
\usepackage{mathtools}
\usepackage{amsthm}
\usepackage{enumitem}
\usepackage{xspace}
\usepackage{xcolor}
\usepackage{listings}
\usepackage{algorithm}
\usepackage{algorithmic}
\usepackage{hyperref}
\usepackage{url}
\usepackage{icml2026}

\icmltitlerunning{ALPS: Adaptive Lineage-Aware Parallel Search for LLM Agents}

\newcommand{\alps}{ALPS\xspace}
\newcommand{\rails}{RAILS\xspace}
\newcommand{\keep}{\texttt{KEEP\_EXPERIMENT}\xspace}
\newcommand{\todo}[1]{\textcolor{red!70!black}{[TODO: #1]}}
\newcommand{\tbd}{\textcolor{red!70!black}{TBD}}
\newcommand{\lowerbetter}{$\downarrow$}
\newcommand{\higherbetter}{$\uparrow$}
\newcommand{\eval}{\mathrm{eval}}
\newcommand{\HEAD}{\textsc{head}\xspace}
\newcommand{\head}{\mathrm{HEAD}}

\lstset{
  basicstyle=\ttfamily\footnotesize,
  breaklines=true,
  frame=single,
  framesep=3pt,
  xleftmargin=0pt,
  xrightmargin=0pt,
  columns=fullflexible,
  keepspaces=true,
  aboveskip=5pt,
  belowskip=5pt
}

\begin{document}

\twocolumn[
\icmltitle{ALPS: Adaptive Lineage-Aware Parallel Search for LLM-Driven Optimization}

\begin{icmlauthorlist}
\icmlauthor{Anonymous Authors}{anon}
\end{icmlauthorlist}
\icmlaffiliation{anon}{Anonymous Institution}
\icmlcorrespondingauthor{Anonymous Author}{anon.email@domain.com}
\icmlkeywords{LLM agents, black-box optimization, automated research, parallel search, MCTS, hyperparameter tuning, data mixture}
\vskip 0.3in
]

\printAffiliationsAndNotice{}

\begin{abstract}
Language-model agents are increasingly used as black-box optimization policies: they propose a code or configuration edit, evaluate it with a fixed harness, observe a scalar metric, and iterate. Serial agent loops use feedback efficiently but under-utilize parallel hardware, whereas naive parallel loops improve utilization but evaluate many candidates against stale or non-composable baselines. We formulate this tension as a scheduling problem over an asynchronous experiment lineage. We propose \alps (Adaptive Lineage-Aware Parallel Search), a scheduler that combines a lineage selector, an operator-level bandit, a commit-hazard controller for adaptive parallelism, and a noise-aware promotion gate. \alps treats the LLM as a candidate proposer and assigns all stateful decisions---dispatch, validation, commit, and rebase---to the scheduler. We evaluate \alps on two tasks: a small-GPT autoresearch benchmark and a Qwen3 supervised fine-tuning data-mixture search. Across both tasks, we compare three policies---serial, naive parallel, and \alps---under matched wall-clock budgets. Preliminary results suggest that lineage-aware scheduling can recover cumulative-improvement behavior while retaining the throughput advantages of parallel execution.
\end{abstract}

\section{Introduction}
\label{sec:intro}

A growing body of work treats large language models (LLMs) not only as predictors, but also as optimization policies that generate candidate solutions, observe feedback, and refine subsequent proposals. This pattern appears in automated research and code-search systems \citep{agentlab,karpathy_autoresearch,ai_scientist,aiscientistv2,aide,alphaevolve,funsearch}, prompt optimization \citep{yang2023opro,promptagent,protegi,evoprompt,promptbreeder}, compound LM-system tuning \citep{textgrad,dspy}, and reasoning-and-acting agents \citep{react,reflexion,yao2023tree,rap,lats}. We use \emph{agent-led search} to denote this family of procedures. In the setting considered here, each decision produces a runnable edit---for example, a hyperparameter change, a data-mixture vector, or a code patch---which is evaluated by a fixed harness and scored by a noisy scalar objective.

The practical deployment regime for such loops is increasingly parallel. A user may have access to an eight-GPU node or a shared cluster allocation, while each evaluation is slow enough that a purely serial loop leaves substantial accelerator time unused. However, serial agent loops have algorithmic advantages that are easy to lose. Because the agent observes the most recent result before proposing the next candidate, each decision is conditioned on fresh feedback. Because the repository baseline advances after each accepted edit, later candidates can compose previous improvements. And because the experiment record is a single trajectory, the agent's reasoning context remains coherent. These properties help explain why simple serial systems such as \texttt{autoresearch} \citep{karpathy_autoresearch} can be effective despite low hardware utilization.

Naive manager-worker parallelism improves utilization but weakens these properties. A manager dispatches $K$ candidates from the same history, workers evaluate them independently, and the system reacts only after a batch or partial batch completes. Within such a batch, improvements cannot compound: every candidate is evaluated against the same baseline, even if one of its siblings would have changed what the next candidate should try. In addition, many evaluations may be launched from a baseline that becomes obsolete before they finish. The manager therefore receives more measurements, but those measurements are less useful for constructing a coherent cumulative trajectory. The result is a tension between \emph{freshness}, which favors serial interaction, and \emph{throughput}, which favors parallel execution.

This paper argues that the tension is not primarily a failure of LLM proposal quality. Rather, it is a missing runtime decision problem. Once multiple LLM-proposed experiments run asynchronously, the system must decide which branch of the source lineage to expand, which edit axis to explore, how many workers to launch speculatively, when an observed improvement is statistically strong enough to promote, and how to handle candidates evaluated on stale baselines. Prompting the LLM manager to ``coordinate'' workers is not enough: the runtime needs an explicit scheduler with access to lineage state, in-flight evaluations, metric noise, and promotion rules.

We propose \alps, Adaptive Lineage-Aware Parallel Search, a scheduling policy for agent-led optimization under parallel evaluation. \alps separates two roles. The LLM remains the generative proposer: it suggests structured edits and provides rationales. The scheduler makes the stateful decisions: where to branch, which operator to sample, how much speculative parallelism to allow, and whether a completed candidate should advance the shared lineage. Concretely, \alps represents the experiment process as asynchronous expansion of a git-backed commit tree. It combines a UCT-style selector over lineage nodes \citep{uct}, an operator-level bandit over edit axes, a commit-hazard controller that adapts queue depth to the probability that \HEAD will soon change, diversity-aware batch construction inspired by batch Bayesian optimization \citep{thompson_batch_bo,local_penalization}, stale-candidate rebase validation, and a noise-aware promotion gate.

The central idea is that parallelism should be adaptive rather than fixed-width. When the scheduler estimates that no in-flight evaluation is likely to advance \HEAD, it can safely fill the worker pool. When an in-flight candidate is likely to commit, launching many additional candidates from the current \HEAD may waste work, because those candidates will soon be evaluated against an obsolete baseline. \alps therefore trades off utilization and freshness explicitly through a commit-hazard estimate. This makes the scheduler different from both serial search, which always uses the freshest state but only one worker, and naive parallel search, which always fills all workers but ignores the changing lineage.

We evaluate \alps on two agent-led optimization tasks that stress different kinds of compositionality. The first is a small-GPT autoresearch task in which the agent edits training code and hyperparameters. The second is a Qwen3 supervised fine-tuning data-mixture task in which the agent edits a constrained mixture over instruction-following domains. Across both tasks, we compare three scheduling policies under matched wall-clock budgets: serial search, naive parallel search, and \alps. This $3 \times 2$ comparison isolates the core tradeoff between information freshness and hardware utilization, and tests whether explicit lineage-aware scheduling can recover cumulative improvement while still using parallel workers.

Our contributions are as follows:
\begin{enumerate}[leftmargin=*,itemsep=2pt]
  \item \textbf{A scheduling formulation for agent-led parallel search.}
  We formalize the gap between serial and naive-parallel agent loops as a decision problem over a mutable experiment lineage, an operator distribution, an in-flight evaluation set, an adaptive speculative queue depth, and a statistical promotion rule.

  \item \textbf{The \alps scheduler.}
  We instantiate this formulation with lineage-aware branch selection, operator-level exploration, commit-hazard adaptive parallelism, diversity-aware batch construction, stale-candidate rebase validation, and a noise-aware commit gate. The design treats the LLM as a proposer and the runtime as the stateful scheduler.

  \item \textbf{A two-task empirical study.}
  We evaluate serial search, naive parallel search, and \alps on small-GPT autoresearch and Qwen3 SFT data-mixture optimization under matched wall-clock budgets. The comparison shows how lineage-aware scheduling changes the tradeoff between evaluation throughput, feedback freshness, and cumulative composition.
\end{enumerate}
\section{Background and Related Work}
\label{sec:related}

\paragraph{LLMs as optimizers and programmatic agents.}
OPRO casts optimization as an in-context procedure in which an LLM is prompted with previous solutions and scores \citep{yang2023opro}. PromptAgent extends this idea with strategic planning for prompt optimization \citep{promptagent}. Related prompt-optimization methods such as APO/ProTeGi, EvoPrompt, and PromptBreeder use textual-gradient, evolutionary, or population-based mechanisms to improve discrete prompts under validation feedback \citep{protegi,evoprompt,promptbreeder}. TextGrad and DSPy optimize compound LM systems by propagating textual feedback or compiling declarative programs into improved pipelines \citep{textgrad,dspy}. ReAct, Reflexion, Tree-of-Thoughts, RAP, and LATS show that language agents can interleave reasoning, acting, memory, and explicit search over action or reasoning trajectories \citep{react,reflexion,yao2023tree,rap,lats}. These methods primarily improve the proposer, memory, program representation, or reasoning trace. \alps is complementary: it asks how an external runtime should schedule, validate, and compose expensive evaluations produced by any such proposer.

\paragraph{Agentic code search and automated research.}
A line of recent work treats LLMs as code-generating search operators inside an evaluator-driven loop. FunSearch pairs an LLM with an evolutionary program-search procedure and a deterministic evaluator to discover mathematical constructions \citep{funsearch}. AIDE frames machine-learning engineering as tree search over code solutions, where LLM-generated patches create child nodes that are evaluated and reused \citep{aide}. The AI Scientist and AI Scientist-v2 broaden this pattern to autonomous scientific workflows, with v2 explicitly using progressive agentic tree search and parallel experiment execution \citep{ai_scientist,aiscientistv2}. AlphaEvolve similarly uses LLM-generated code edits, evolutionary selection, and programmatic evaluators for scientific and systems optimization \citep{alphaevolve}. Concurrent benchmarks such as MLE-bench and RE-Bench evaluate ML-engineering and AI-R\&D agents under realistic experiment budgets \citep{mlebench,rebench}. These systems demonstrate the importance of evaluator-driven iteration, but they mainly study agent capability, search operators, or full research workflows; \alps instead focuses on the runtime scheduling problem that arises when expensive LLM-proposed evaluations run asynchronously and may compose through a changing source lineage.

\paragraph{Tree search.}
Upper-Confidence Trees (UCT) use confidence bonuses to balance branch exploitation and exploration in Monte-Carlo planning \citep{uct}. Agent-led optimization has an analogous tree structure because accepted edits create new source states, and subsequent candidates may be applied either to the current \HEAD or to earlier branches. The distinction is that rollouts are expensive wall-clock evaluations, rewards are noisy, and the candidate generator is a language model rather than a simple simulator policy. \alps therefore uses UCT only for lineage selection and delegates candidate generation to the LLM proposer.

\paragraph{Parallel hyperparameter optimization.}
Hyperband and ASHA achieve parallelism by allocating resources adaptively and terminating weak configurations early \citep{hyperband,asha}; BOHB combines Hyperband-style resource allocation with model-based Bayesian optimization \citep{bohb}. Production and open-source systems such as Google Vizier and Optuna further emphasize that black-box optimization at scale requires not only an acquisition policy but also robust trial management, distributed execution, and fault tolerance \citep{vizier,optuna}. Population-based training (PBT) maintains a population of models and alternates exploitation with perturbation \citep{pbt}. These methods motivate two aspects of \alps: explicit resource scheduling and the treatment of good stale candidates as objects that may be rebased and re-evaluated. However, agent-led search differs from conventional HPO in that candidates are structured source edits with natural-language rationales, and accepting a candidate changes the semantic baseline on which future candidates are meaningful.

\paragraph{Batch Bayesian optimization.}
Parallel Bayesian optimization constructs batches that jointly trade off value and diversity. Parallel Thompson sampling samples multiple functions from the posterior, whereas local penalization discourages selecting nearby points within the same batch \citep{thompson_batch_bo,local_penalization}. \alps adopts the same principle at the level of edit operators and rationales: batch members should not be minor variants of the same hypothesis unless the scheduler is deliberately validating noise. More general Bayesian-optimization treatments also emphasize noisy expensive evaluations, asynchronous or parallel acquisition, and the distinction between choosing a batch and assimilating partial results \citep{frazier2018tutorial,parallel_ts,qei}. These assumptions are close to our wall-clock setting, but continuous-space BO usually treats candidates as independent design points rather than edits whose value depends on a mutable code lineage.

\paragraph{Autonomous research and data-mixture tuning.}
Agent Laboratory studies broader research-assistant workflows \citep{agentlab}, while \texttt{autoresearch} isolates the inner optimization loop of proposing, evaluating, and keeping or reverting edits \citep{karpathy_autoresearch}. The latter is a useful testbed because it exposes the scheduling problem without requiring a full research pipeline. Our second pilot uses Qwen3-0.6B \citep{qwen3} and T\"ulu-3 instruction-following data \citep{tulu3}; the availability of labeled data sources makes data-mixture search interpretable and compositional, and thus well matched to structured agent-led optimization.

\section{Preliminaries}
\label{sec:problem}

We formalize agent-led parallel search as a scheduling problem and introduce the small set of state quantities the policy in Section~\ref{sec:alps} consumes.

\paragraph{Search target.}
A search task supplies a runnable codebase ${\cal C}_0$ and a harness $f(\cdot)$ that takes a code state and returns a noisy scalar score. Given a working baseline $b_t$ (a particular commit of the codebase visible at time $t$), a \emph{candidate} is a structured edit $u_t$ that yields a new state $s_t = u_t(b_t)$. The harness produces
\begin{equation}
  y_t = f(s_t) + \epsilon_t, \qquad \epsilon_t \sim \mathcal{D}(0, \sigma^2),
\end{equation}
where lower is better. We assume the structured-edit form gives every candidate a discrete \emph{operator} label $o_t$ (which axis the edit targets, e.g.\ \texttt{ASPECT\_RATIO} or the data-mix \texttt{DATA\_MIX} dict) and a value $v_t$. The proposer is an LLM policy $\pi_\theta$ that conditions on the history $H_t = \{(b_i, o_i, v_i, r_i, y_i)\}_{i<t}$, where $r_i$ is the rationale the LLM recorded for $u_i$. The objective is to minimize the best-so-far score under wall-clock budget $W$ with $K$ parallel workers.

\paragraph{Cumulative lineage.}
Each successful $y$ either is discarded or is \emph{committed} into a shared baseline by re-applying its edit to the current \HEAD and advancing \HEAD by one git commit. A run produces a chain (or tree) of commits $b_0 \to b_1 \to \cdots \to b_t$ which we call the \emph{lineage}. Subsequent dispatches clone from $b_t$, so accepted improvements compound. This is the primitive that makes Karpathy-style serial agents effective and that naive parallel dispatch lacks; making it work in parallel is the technical contribution of this paper.

\paragraph{The five scheduling decisions.}
At any wake-up, a parallel agent-led system implicitly makes five decisions:

\begin{itemize}[leftmargin=*,nosep]
\item Which lineage \emph{node} $b$ to expand from.
\item Which \emph{operator} $o$ to apply (axis-level exploration vs.\ exploitation).
\item How many workers to keep \emph{speculatively} dispatched against the current \HEAD vs.\ to hold for in-flight outcomes.
\item When to \emph{commit} a finished candidate back into the lineage — a noisy decision whose downside is invalidating concurrent in-flight evaluations.
\item When to \emph{validate} a stale-baseline candidate via a rebased re-evaluation rather than commit or discard.
\end{itemize}

A serial loop answers all five trivially (always \HEAD; whatever the agent proposes; one worker; advance on any improvement; staleness does not arise). A naive parallel loop also answers all five trivially (always \HEAD; whatever the manager proposes; $K$ workers always speculative; commit at batch barriers; no validation). Both are degenerate scheduling policies. The interesting question is what a non-degenerate policy looks like.

\paragraph{Commit hazard.}
The bridge between in-flight state and queue depth is the \emph{commit hazard} $\rho_t$, the probability that some running evaluation will advance \HEAD when it finishes:
\begin{equation}
  \rho_t \;=\; 1 - \prod_{e \in \mathcal{R}_t} \bigl(1 - \mathbb{P}\!\left[\,e \text{ commits over HEAD}\,\right]\bigr),
  \label{eq:hazard}
\end{equation}
where $\mathcal{R}_t$ is the set of in-flight evaluations. The per-evaluation commit probability is computed from the operator posterior (Section~\ref{sec:alps:operators}); when posteriors are uninformative (early in a run), it falls back to a fixed prior. $\rho_t$ summarizes the information-vs-speculation tradeoff: a large $\rho_t$ says today's \HEAD will likely be wrong soon, so speculative work launched against it is at risk; a small $\rho_t$ says \HEAD is stable and the safe move is to fill the worker pool.

\paragraph{Per-task noise estimate.}
The commit decision is statistical because $\epsilon_t$ is non-trivial. We measure $\hat\sigma$ from a small number of \emph{calibration repeats}: identical no-edit runs of the harness on the current \HEAD, with varied PRNG seeds. This is task-specific (each harness has its own seed-to-seed spread) and is a one-time cost per source state. In our autoresearch experiments three calibration repeats give $\hat\sigma \approx 3 \times 10^{-3}$ on a 10-min training budget, an order of magnitude larger than what we initially assumed and large enough that without an explicit gate every candidate would look like a winner.

\section{ALPS: The Scheduling Policy}
\label{sec:alps}

\alps is an event-driven scheduler for agent-led optimization. It separates two roles: the LLM proposes candidate edits, while the scheduler decides where to branch, which edit axis to explore, how much parallel speculation to allow, and when a completed candidate should change the shared baseline. The scheduler maintains three pieces of state: a git-backed lineage tree, an operator posterior over edit axes, and the set of running evaluations.

\subsection{Scheduler state}

Each lineage node $v$ corresponds to a committed source state and stores its best observed metric $\hat y_v$, the number of evaluations expanded from it $n_v$, and the edit history leading to it. Each candidate also has an operator label $o \in \mathcal{O}$, such as an architecture edit, optimizer edit, schedule edit, validation repeat, or data-mixture edit. For each operator, \alps maintains an online posterior over expected improvement,
\[
  \Delta_o \sim \mathcal{N}(\mu_o,\sigma_o^2),
\]
updated from completed evaluations. This state is deliberately lightweight: it uses telemetry that the runtime already records rather than requiring a learned surrogate model.

\subsection{Selecting a lineage and an operator}

At each dispatch step, \alps first selects a lineage node using a UCT-style score:
\begin{equation}
  S_{\mathrm{lineage}}(v)
  =
  -\hat y_v
  +
  \beta \sqrt{\frac{\log N}{1+n_v}}
  -
  \lambda_{\mathrm{stale}}\mathrm{stale}(v),
  \label{eq:uct}
\end{equation}
where lower metric is better, $N$ is the total number of completed evaluations, and $\mathrm{stale}(v)$ penalizes branches that are far from the current \HEAD and therefore harder to rebase.

It then selects an operator using an uncertainty-aware bandit score:
\begin{equation}
  S_{\mathrm{op}}(o)
  =
  \mu_o
  +
  \alpha\sigma_o
  +
  \gamma\mathrm{cov}(o)
  -
  \eta\mathrm{fail}(o)
  -
  \rho\mathrm{cost}(o).
  \label{eq:op}
\end{equation}
The first term exploits operators that have produced improvements; the second and third terms encourage uncertain or under-covered axes; the last two terms penalize invalid edits and expensive evaluations.

\subsection{Adapting parallelism by commit hazard}

The key difference between \alps and fixed-width parallel search is that \alps does not always keep all workers busy. It estimates the probability that at least one running evaluation will advance the current \HEAD:
\begin{equation}
  \rho_t
  =
  1 -
  \prod_{e \in \mathcal{R}_t}
  \left(
  1 -
  \Pr[e \text{ commits over } \HEAD]
  \right),
  \label{eq:hazard}
\end{equation}
where $\mathcal{R}_t$ is the set of in-flight evaluations. This commit hazard controls the target queue depth:
\begin{equation}
  Q_t
  =
  Q_{\min}
  +
  (Q_{\max}-Q_{\min})(1-\rho_t)^\zeta .
  \label{eq:queue}
\end{equation}
When $\rho_t$ is high, \HEAD is likely to change soon, so launching more speculative work from the current baseline is risky. When $\rho_t$ is low, the baseline appears stable and \alps fills more workers.

\subsection{Batch diversity}

When multiple workers are available, \alps constructs a batch greedily. A candidate $c$ has a standalone acquisition value $A(c)$ from the selected lineage and operator scores. Within a batch $B$, the value is penalized by similarity to already selected candidates:
\begin{equation}
  A(c \mid B)
  =
  A(c)
  -
  \delta
  \sum_{c' \in B}
  \mathrm{sim}(c,c').
  \label{eq:batch}
\end{equation}
The similarity function can be simple: candidates with the same operator and value are near-duplicates; candidates with the same operator but different values are partially similar; candidates on different operators are dissimilar. This prevents a parallel batch from collapsing into minor variants of the first promising hypothesis.

\subsection{Commit and rebase}

When an evaluation finishes, \alps decides whether it should change the shared lineage. Let $y_e$ be the candidate score, $y_{\HEAD}$ the current head score, and $\hat\sigma$ the task-level noise estimate from calibration repeats. \alps commits only if
\begin{equation}
  y_{\HEAD} - y_e
  >
  \tau_t\hat\sigma
  +
  c_{\mathrm{stale}}N_{\mathrm{stale}},
  \label{eq:gate}
\end{equation}
where $N_{\mathrm{stale}}$ is the number of running evaluations that would become stale if the commit fires. The threshold $\tau_t$ can decay over the run, making the scheduler conservative early and more permissive near the deadline.

If the candidate was evaluated on the current \HEAD, the edit is committed immediately. If it was evaluated on an older baseline, \alps does not transfer the metric directly. Instead, it re-applies the same edit to the current \HEAD and schedules a rebase-validation evaluation. This keeps the lineage cumulative while avoiding invalid commits from stale baselines.

\subsection{Algorithm}

\begin{algorithm}[!ht]
\caption{\alps scheduler}
\label{alg:alps}
\begin{algorithmic}[1]
\STATE Initialize lineage tree $G$ at $b_0$, operator posteriors $\{\mu_o,\sigma_o\}$, and running set $\mathcal{R}\leftarrow\emptyset$.
\WHILE{wall-clock budget remains}
  \STATE Collect newly finished evaluations.
  \FOR{each finished evaluation $e$}
    \STATE Update lineage statistics and operator posterior.
    \IF{$e$ passes the commit gate in Eq.~\ref{eq:gate}}
      \IF{$e$ was evaluated on current \HEAD}
        \STATE Apply its edit, advance \HEAD, and log the commit.
      \ELSE
        \STATE Enqueue a rebase-validation task on current \HEAD.
      \ENDIF
    \ENDIF
  \ENDFOR

  \STATE Compute commit hazard $\rho_t$ by Eq.~\ref{eq:hazard}.
  \STATE Set target queue depth $Q_t$ by Eq.~\ref{eq:queue}.

  \WHILE{$|\mathcal{R}| < Q_t$ and workers are available}
    \STATE Select lineage node $v$ by Eq.~\ref{eq:uct}.
    \STATE Select operator $o$ by Eq.~\ref{eq:op}.
    \STATE Ask the LLM proposer for candidates constrained to $(v,o)$.
    \STATE Select the candidate with highest diversity-adjusted acquisition in Eq.~\ref{eq:batch}.
    \STATE Dispatch it to a worker and add it to $\mathcal{R}$.
  \ENDWHILE
\ENDWHILE
\STATE \textbf{return} best committed source state.
\end{algorithmic}
\end{algorithm}


\section{System Implementation}
\label{sec:system}
\begin{figure}
    \centering
    \includegraphics[width=1\linewidth]{image.png}
    \caption{System Design}
    \label{fig:placeholder}
\end{figure}
\alps requires an execution layer that can run LLM-proposed candidates in parallel while preserving a canonical source lineage. We implement a lightweight task runtime for this purpose. The runtime is not part of the algorithmic contribution; it provides the mechanisms needed by the scheduler: isolated execution, structured edits, metric logging, commit validation, and rebase evaluation.

\paragraph{Candidate representation.}
Each candidate is represented as
\[
  c = (b, o, u, r),
\]
where $b$ is the base commit, $o$ is the operator label, $u$ is a structured edit, and $r$ is the LLM-provided rationale. The operator label feeds the bandit in Eq.~\ref{eq:op}; the edit is applied to the selected base commit; and the rationale is stored for auditability and optional similarity computation.

\paragraph{Isolated evaluation.}
Before execution, the runtime clones the selected base commit into a per-evaluation sandbox, applies the candidate edit, and runs the task harness. The harness returns a scalar metric and metadata. This isolation prevents parallel evaluations from racing on a shared source tree and ensures that every metric can be traced to an exact source state.

\paragraph{Lineage update.}
The canonical state is a git lineage. When a candidate passes the commit gate in Eq.~\ref{eq:gate}, the runtime reapplies its edit to the current \HEAD, validates the patch, commits the resulting source state, and records the event. Future candidates can then build on this updated \HEAD, allowing improvements to compose across evaluations.

\paragraph{Rebase validation.}
If a candidate was evaluated on an older baseline, its metric is not transferred directly. Instead, the runtime schedules a rebase-validation evaluation by applying the same edit to the current \HEAD and obtaining a fresh measurement. This converts stale work into an explicit validation step rather than silently discarding it or incorrectly committing it.

\paragraph{Task interface.}
A task only needs to provide a runnable codebase, a harness that emits a scalar metric, and a prompt describing the valid edit axes. The same scheduler can therefore be used for architecture/code search and SFT data-mixture search.


\section{Experiments}
\label{sec:experiments}

We evaluate three scheduling policies on two agent-led optimization tasks under matched wall-clock budgets. The policies are: \textbf{Serial}, a single-worker loop that always evaluates from the latest \HEAD; \textbf{Naive parallel}, an eight-worker batch baseline that evaluates candidates independently from the same baseline; and \textbf{\alps}, our lineage-aware scheduler. The two tasks are a small-GPT autoresearch code-search task and a Qwen3 SFT data-mixture search task.

\subsection{Tasks and metrics}

\paragraph{Autoresearch.}
The first task is a small-GPT training-code search problem derived from \texttt{autoresearch} \citep{karpathy_autoresearch}. The agent edits training hyperparameters or code-level choices in \texttt{train.py}. Each evaluation trains for a fixed budget, and the metric is held-out validation bits-per-byte, where lower is better.

\paragraph{Qwen3 SFT data-mixture search.}
The second task tunes the supervised fine-tuning data mixture for Qwen3-0.6B \citep{qwen3} using T\"ulu-3 instruction-following data \citep{tulu3}. The agent edits a constrained \texttt{DATA\_MIX} dictionary over five domains: math, code, chat, instruction following, and reasoning. Each evaluation performs a short SFT run and reports balanced held-out validation loss across domains.

For both tasks, the primary metric is the best validated score reached before the wall-clock deadline. We also report the number of completed evaluations as a measure of throughput.

\subsection{Main results}

\begin{table*}[t]
\centering
\small
\begin{tabular}{llccccp{4.8cm}}
\toprule
Task & Policy & Workers & Wall-clock & Evals & Best \lowerbetter & Observation \\
\midrule
Autoresearch
& Serial
& 1
& 2 h
& 4
& 0.9750
& Best current score; benefits from fresh feedback and immediate composition. \\

Autoresearch
& Naive parallel\textsuperscript{$\dagger$}
& 8
& 2 h
& 44
& 0.9449
& More evaluations, but little cumulative improvement. \\

Autoresearch
& \alps
& 8
& 2 h
& 50
& 0.9369
& Highest throughput; one late lineage commit composes multiple discovered edits. \\

\midrule
Qwen3 SFT mix
& Serial
& 1
& 4 h
& 9
& 1.10553
& Single agent commits each accepted edit; throughput limited but lineage advances 2 times. \\

Qwen3 SFT mix
& Naive parallel
& 8
& 4 h
& 49
& 1.10168
& Manager retains \keep at wave-end but did not fire it; advances were captured passively from a concurrent run that wrote the canonical source. \\

Qwen3 SFT mix
& \alps
& 8
& 4 h
& 37
& 1.10089
& Explicit gate auto-fires KEEP; manager autonomously explores off-axis operators (LR, weight\_decay, warmup) once the framework declares them as search axes. \\

\bottomrule
\end{tabular}
\caption{
Main $3 \times 2$ comparison. We compare three scheduling policies on two tasks under matched wall-clock budgets. Lower metric is better for both tasks. \textsuperscript{$\dagger$}
}
\label{tab:main_results}
\end{table*}

Table~\ref{tab:main_results} shows the current pilot results. On                  
  autoresearch, serial search remains a strong baseline despite                                     
  completing only eight evaluations, suggesting that fresh feedback                                 
  and immediate composition are valuable. Naive parallel search
  completes 5\(\times\) more evaluations but produces little improvement,                           
  consistent with the failure mode that candidates do not build on one                              
  another. \alps completes 6\(\times\) more evaluations than serial and
  commits five lineage steps inside the run, an independently                                       
  discovered chain of edits: \texttt{ASPECT\_RATIO=96},                                             
  \texttt{TOTAL\_BATCH\_SIZE=$2^{18}$}, \texttt{WEIGHT\_DECAY=0.15},                                
  \texttt{SCALAR\_LR=0.4} (this one auto-rebased after a stale-baseline                             
  gate decision), and \texttt{WARMDOWN\_RATIO=0.9}. This demonstrates                               
  that lineage-aware scheduling can recover the cumulative-improvement                              
  behavior of the serial loop under parallel execution.  

On the Qwen3 SFT data-mixture task, all three policies are now collected. \alps completes 37 evaluations in four hours on eight GPUs and reaches a best balanced val\_loss of 1.10089, against 1.10168 for naive parallel (49 evaluations) and 1.10553 for serial (9 evaluations). The empirical noise floor estimated from in-run calibration repeats is $\hat\sigma \approx 1.9 \times 10^{-5}$ (3-sample) to $5.3 \times 10^{-5}$ (5-sample), so the gap between \alps and serial is more than an order of magnitude above any reasonable noise threshold. The \alps run also illustrates a behavioral effect that did not arise in the autoresearch task: once the framework declared \texttt{LR}, \texttt{weight\_decay}, and \texttt{warmup} as registered operators in addition to \texttt{DATA\_MIX}, the manager autonomously proposed edits along those axes and the gate accepted them. The headline improvement on this task is therefore a mixture of data-mix tilt (math$\approx$0.30, reasoning$\approx$0.27) and a learning-rate increase (\texttt{LR}=3e-5 vs the default 2e-5), demonstrating that the operator-level bandit registers exactly the axes the framework exposes.

\subsection{Takeaways}

The pilot results support two preliminary takeaways. First, evaluation count alone is not enough: naive parallelism can improve throughput without improving the best committed state. Second, lineage matters when edits are compositional. \alps is designed to preserve this compositional structure while still using parallel workers. The remaining experimental work is to complete the $3 \times 2$ comparison, including the SFT serial and naive-parallel baselines and a rerun of the autoresearch naive-parallel baseline with isolated sandboxes.


\section{Future Work}
\label{sec:future}

ALPS suggests several directions for future work. First, broader evaluations on prompt optimization, architecture search, systems tuning, and longer autonomous-research workflows would test whether lineage-aware scheduling generalizes beyond the two tasks studied here. Second, targeted ablations can isolate the contribution of each scheduler component, including commit-hazard control, operator-bandit selection, diversity-aware batching, stale-candidate rebase validation, and noise-aware promotion. Third, the current hand-designed scheduling rules could be replaced or augmented with learned policies trained from accumulated optimization traces. Finally, extending ALPS to multi-objective settings, such as accuracy--cost--latency tradeoffs or safety-constrained optimization, would make the framework more applicable to realistic agent-led optimization.

\section{Conclusion}
\label{sec:conclusion}

Agent-led optimization faces a structural tradeoff between the information efficiency of serial search and the hardware utilization of naive parallel search. We argued that this tradeoff should be addressed as a scheduling problem over a commit lineage. \alps implements this view through lineage selection, operator-level exploration, hazard-adaptive parallelism, diversity-aware batching, stale-winner rebase, and noise-aware promotion. We also provides the runtime substrate needed to execute these decisions. Preliminary manual-\alps pilots suggest that lineage-aware scheduling can produce cumulative improvements under parallel execution, while also revealing that prompt-level scheduling is too conservative and axis-fixated.

\section*{LLM Usage Disclosure}
The writing of the article and the implementation of the system were assisted by LLM. System design, and analysis were conducted by the authors.

\bibliographystyle{icml2026}
\bibliography{references}


\end{document}
